% discretization.tex
%
% Purpose:
%   Summarize the Nyström, Kress, MFS/proxy, and implementation components used
%   by the report-level formulation.
%
% Main algorithm:
%   Describes how smooth boundary integral operators and singular kernels are
%   discretized, and how quasiperiodic Green functions are approximated.
%
% Based on:
%   notes/implementation/kress_quadrature.md, notes/implementation/kress_scheme.md,
%   notes/implementation/mfs2d.md, and the package layout in the repository.
%
% Main changes:
%   The material is shortened to a high-level implementation description with
%   the key formulas preserved.
%
% Numerical goal:
%   Identify the discretization choices and package roles that should be varied
%   during convergence checks.

\section{Discretization and Implementation Components}

\subsection{Boundary quadrature}

The inclusion boundaries are assumed smooth and are discretized by a periodic
Nystr\"om method.  For a \(2\pi\)-periodic boundary parameter \(t\), the nodes are
\[
  t_j=\frac{2\pi j}{N},
  \qquad
  j=0,\ldots,N-1,
\]
with even \(N\).  Smooth integral blocks use the periodic trapezoid rule, which
is spectrally accurate for smooth periodic data.

Logarithmic singularities are handled by Kress product quadrature
\cite{Kress1991,Kress2014}.  The kernel is split as
\[
  k(t,s)
  =
  \phi(t,s)\log\!\left(4\sin^2\frac{t-s}{2}\right)
  +
  \psi(t,s),
\]
where \(\phi\) and \(\psi\) are smooth.  The Kress weights for the singular
factor are
\[
  R_j^{(N/2)}(t)
  =
  -\frac{4\pi}{N}
  \left[
  \sum_{n=1}^{N/2-1}\frac{1}{n}\cos(n(t_j-t))
  +
  \frac{1}{N}\cos\!\left(\frac{N}{2}(t_j-t)\right)
  \right].
\]
The corresponding Nystr\"om matrix has entries of the form
\[
  A_{ij}=R_{ij}\phi(t_i,t_j)+h\psi(t_i,t_j),
  \qquad
  h=\frac{2\pi}{N}.
\]
For the Muller matrix, hypersingular differences such as
\[
  T^{(k_e)}-T^{(k_i)}
\]
are treated in a cancellation-aware Kress form.  The leading singular parts
cancel to a lower-order structure, but the implementation still needs the
correct split and diagonal treatment.

\subsection{Quasiperiodic Green functions and proxy/MFS periodization}

The quasiperiodic Green function is one of the numerically delicate parts of the
method.  In a one-direction periodic open waveguide, the Green function has a
Rayleigh expansion in the open direction, for example
\[
  G_{\mathrm{QP}}^{(k)}((x,y),(x_s,y_s))
  =
  \frac{\ii}{2d}
  \sum_{m\in\mathbb Z}
  \frac{1}{\gamma_m}
  \ee^{\ii\beta_m(y-y_s)}
  \ee^{\ii\gamma_m|x-x_s|},
\]
away from thresholds.  Direct lattice sums may be slowly convergent or singular
near exceptional parameters.

The codebase uses MFS/proxy ideas to approximate quasiperiodic corrections.  In
a typical proxy representation,
\[
  G_{\mathrm{QP}}^{(k)}(x,y)
  \approx
  G^{(k)}(x,y)
  +
  \sum_{\ell=1}^{P}q_\ell(y)\,G^{(k)}(x,z_\ell),
\]
where \(G^{(k)}\) is the free-space singular Green function and the proxy
sources \(z_\ell\) lie outside the physical cell.  The proxy strengths enforce
quasiperiodic value and derivative discrepancies on the cell boundary.  The
free-space singular part is retained; the proxy correction is a homogeneous
field inside the cell.

This periodization strategy is flexible and useful for the present experiments,
but its parameters require convergence checks.  It should be compared with, and
understood alongside, alternative quasiperiodic integral representations such as
those discussed by Barnett and Greengard \cite{BarnettGreengard2010}.

\subsection{Far-field extractors}

The matrices \(F_\pm\) map the center boundary density to outgoing Rayleigh
coefficients at the ports:
\[
  s^- = F_-\eta,
  \qquad
  s^+ = F_+\eta .
\]
At the continuous level, the right-going coefficient functions come from the
Rayleigh expansion of \(G_{\mathrm{QP}}\).  For instance, for a source point
\(z=(z_x,z_y)\) and right reference wall \(X_R\),
\[
  g_m^R(z)
  =
  \frac{\ii}{2\sqrt d\,\gamma_m}
  \exp[-\ii\beta_m z_y-\ii\gamma_m(z_x-X_R)] .
\]
The double-layer source-normal contribution has coefficient
\[
  -\ii(\gamma_m\nu_x+\beta_m\nu_y)g_m^R(z),
\]
and the single-layer contribution is proportional to \(-g_m^R(z)\) under the
current \(\eta=[\tau;-\sigma]\) convention.  The left extractor is analogous
with the left-going Rayleigh phase and the corresponding source-normal
derivative formula.

\subsection{Implementation components}

The present code organization separates the numerical roles into packages and
driver scripts:
\[
\begin{array}{ll}
\texttt{+geom} & \text{geometry construction and cached contour data},\\
\texttt{+kernel} & \text{quasiperiodic Green, MFS/proxy, and Kress kernel data},\\
\texttt{+op} & \text{Muller matrix assembly, including }\texttt{construct\_A\_QP},\\
\texttt{+bloch} & \text{Rayleigh channels, scattering matrices, modes, traces, selection},\\
\texttt{+viz} & \text{geometry and field visualization helpers}.
\end{array}
\]
The scattering workflows are represented by scripts such as
\path|scat_ld_lead_in.m|, \path|scat_edc_lead_in.m|, and
\path|scat_edc_ps.m|.  Eigenvalue or transmission-eigenvalue scans are
represented by scripts with the \path|tep_| prefix.  These names are
implementation anchors rather than a fixed API specification.
